\documentclass[12pt]{article}
\usepackage[margin=0.75in]{geometry}

\newcommand{\duedate}{04/01/2026}
\newcommand{\assignment}{Learn-to-Learn without Forgetting: Summary}

\newcommand{\name}{Aksel Kretsinger-Walters, adk2164}
\newcommand{\email}{adk2164@columbia.edu}

\makeatletter
\def\input@path{{../}{../../}{../../../}}
\makeatother
\input{pset_template_CLMM.tex} %% DO NOT CHANGE THIS LINE

\begin{document}
\psetheader %% DO NOT CHANGE THIS LINE

\section*{Learning to Learn without Forgetting by Maximizing Transfer and Minimizing Interference}

\paragraph{Motivation.}
Continual learning over non-stationary distributions remains a major challenge
for neural networks. Prior work has framed the problem primarily through the
lens of catastrophic forgetting (interference) or the stability-plasticity
dilemma. However, these views consider only backward-in-time effects: how
current learning degrades past knowledge.

This paper argues that a complete picture must be \emph{temporally symmetric},
accounting for both backward interference and forward transfer simultaneously.

\paragraph{The Transfer-Interference Trade-off.}
The authors define operational measures of transfer and interference between
two examples $(x_i, y_i)$ and $(x_j, y_j)$ via gradient alignment:

\begin{itemize}
		\item \textbf{Transfer} occurs when
				$\frac{\partial L(x_i, y_i)}{\partial \theta} \cdot
				\frac{\partial L(x_j, y_j)}{\partial \theta} > 0$,
				meaning learning one example improves the other.
		\item \textbf{Interference} occurs when this dot product is negative,
				meaning learning one example harms the other.
\end{itemize}

Weight sharing between examples creates both transfer and interference potential.
The key insight is that we must consider this trade-off in both the forward and
backward temporal directions, not just backward as in the classical
stability-plasticity framing.

\paragraph{Ideal Objective.}
To maximize transfer and minimize interference, the paper proposes the following
idealized objective:
\[
		\theta = \arg\min_{\theta} \;
		\mathbb{E}_{[(x_i, y_i),(x_j, y_j)] \sim D}
		\Bigl[ L(x_i, y_i) + L(x_j, y_j)
    - \alpha \frac{\partial L(x_i, y_i)}{\partial \theta} \cdot
		\frac{\partial L(x_j, y_j)}{\partial \theta} \Bigr].
\]

The gradient dot-product regularizer directly encourages the network to share
parameters where gradient directions align and keep them separate where
gradients conflict.

\paragraph{Meta-Experience Replay (MER).}
Computing second derivatives of the ideal objective is expensive, so the authors
approximate it using first-order meta-learning (Reptile) combined with experience
replay. The resulting algorithm, \textbf{Meta-Experience Replay (MER)}, works
as follows:

\begin{enumerate}
		\item Maintain an experience replay buffer $M$ updated via reservoir sampling.
		\item At each time step, draw $s$ batches of $k{-}1$ replay samples, each
				interleaved with the current example.
		\item Process each batch with inner SGD steps, applying a within-batch
				Reptile meta-update after each batch.
		\item Apply an across-batch Reptile meta-update to the main parameters.
		\item Add the current example to the replay buffer.
\end{enumerate}

This procedure approximately optimizes the transfer-interference objective with
minimal computational overhead beyond standard experience replay.

\paragraph{Key Design Choices.}
Several properties make MER effective:

\begin{itemize}
		\item \textbf{Reservoir sampling} ensures the buffer approximates the
				stationary distribution over all examples seen so far.
		\item \textbf{Prioritizing the current example} by interleaving it with
				every replay batch ensures the model can learn the present while
				consolidating past knowledge.
		\item \textbf{SGD batch size of 1} provides fine-grained second derivative
				information at the example level.
\end{itemize}

\paragraph{Experimental Results.}
MER is evaluated on supervised continual learning benchmarks and reinforcement
learning settings:

\begin{itemize}
		\item \textbf{MNIST Rotations/Permutations (20 tasks):} MER outperforms
				GEM, EWC, and online baselines in retained accuracy, showing superior
				balance between transfer and interference.
		\item \textbf{Many Permutations \& Omniglot:} In more non-stationary
				settings (100 tasks, 50 alphabets), MER's advantage over baselines
				grows substantially, nearly quadrupling baseline performance on Omniglot.
		\item \textbf{Buffer sensitivity:} MER's benefits increase as the buffer
				shrinks, providing $>$10\% gains over GEM at the smallest buffer sizes.
		\item \textbf{Continual RL (Catcher, Flappy Bird):} DQN-MER retains
				knowledge across task switches while standard DQN forgets
				catastrophically.
\end{itemize}

Gradient dot-product analysis confirms that MER produces significantly more
positive gradient alignment than standard experience replay.

\paragraph{Limitations.}
MER requires maintaining an experience replay buffer, which may be infeasible
at very large scales. The Reptile-based meta-learning adds computational
overhead proportional to the number of replay batches. The theoretical
analysis relies on first-order approximations of the ideal objective.

\paragraph{Takeaway.}
This paper reframes continual learning as solving the transfer-interference
trade-off in a temporally symmetric manner. MER provides a practical,
computationally efficient algorithm that combines experience replay with
meta-learning to promote gradient alignment, consistently outperforming
prior continual learning methods across supervised and reinforcement
learning settings.

\end{document}
